\begin{exercice*}
    Dans un jeu de 32 cartes, il y a quatre familles : les trèfles, les cœurs, les carreaux et les piques, et deux couleurs : rouge et noir. On décide de tirer au hasard une carte dans ce jeu.

    \begin{enumerate}
        \item Si l'on considère la couleur des cartes, cite les issues possibles.
        \item Si l'on considère les familles, cite les issues possibles.
        \item Cite un événement impossible.
        \item Quelle est la probabilité de tirer un trèfle ?
        \item Écris à l'aide d'une phrase non négative, l'événement contraire de « La carte choisie est un trèfle. »
        \item Calcule la probabilité de cet événement.
    \end{enumerate}
\end{exercice*}
\begin{corrige}
    %\setcounter{partie}{0} % Pour s'assurer que le compteur de \partie est à zéro dans les corrigés
    %\phantom{rrr}    
    \begin{enumerate}
        \item Les issues possibles sont : rouge ou noir.
        \item Les issues possibles sont : trèfle, cœur, carreau ou pique.
        \item Un événement impossible est : « Tirer une carte bleue. »
        \item La probabilité de tirer un trèfle est de $\dfrac{1}{4}$, car il y a 4 familles de cartes et elles sont équiprobables.
        \item L'événement contraire de « La carte choisie est un trèfle » est : « La carte choisie n'est pas un trèfle. »
        \item La probabilité de cet événement est de $\dfrac{3}{4}$, car il y a 3 familles de cartes qui ne sont pas des trèfles (cœurs, carreaux et piques) sur les 4 familles au total.
    \end{enumerate}
\end{corrige}

